\section{The torsion Picard scheme under a birational map}
We explain how finite flat torsors detect all of the torsion Picard
scheme, including its nonreduced structure. This deduces the Picard form
of the birational-invariance theorem from known finite-torsor invariance;
it makes no claim of a new birational-invariance principle.
Li--Wang's Corollary 4.6 \cite[p.~7]{LiWang} gives the latter through the
Nori fundamental group scheme. We use Moret--Bailly's purity theorem
below to make the particular Picard deduction explicit.
The coherent vanishing in
Theorem~\ref{thm:birational} is the theorem of Ishii--Yoshida
\cite[Corollary 1.4, p.~3]{IshiiYoshida}: they prove it in every
characteristic for a normal source and regular target, without an
excellence assumption. Their Theorem 1.3 treats more general locally
quasi-unmixed targets satisfying $S_3$ and pseudo-rationality in
codimension two. Applying their corollary after localization also
covers a noetherian target without a global finite-dimensional bound.

For a proper birational map $f:X\to Y$ with $Y$ normal, there is
an open $U\subset Y$, with complement of codimension at least two,
over which $f$ is an isomorphism: the inverse rational map extends
over the codimension-one points by properness. Two finite flat torsors
over a normal integral scheme that agree on a dense open agree
everywhere. Their isomorphism scheme is finite, and a rational section
of a finite scheme extends over a normal domain by integral closure
\cite[Proposition 3, p.~432]{Marrama}.

\subsection{Finite subgroups detect a proper group scheme}
For the Picard assertion we may use the full purity theorem.
Moret--Bailly's theorem says that finite locally free torsors on a
regular noetherian scheme extend uniquely across a closed subset
of codimension at least two \cite[Theorem 6]{Marrama}. It permits
noncommutative group schemes. Combining it with the uniqueness
principle proves equivalence of finite flat $G$-torsors under the
proper birational pullback $f$.

\begin{lemma}\label{lem:detect}
Let $k$ be algebraically closed of characteristic $p$, let $P$
be a proper commutative group scheme of finite type over $k$, and
let $Q\subset P$ be a closed subgroup. If every homomorphism
from a finite commutative $k$-group scheme to $P$ factors through
$Q$, then $Q=P$.
\end{lemma}
\begin{proof}
The reduced subgroup is smooth; its identity component $B$ is
an abelian variety. For $\ell\ne p$, the union of $B[\ell^n]$
is Zariski dense in $B$. Indeed its closure is a subgroup, and
a subgroup of smaller dimension has too few $\ell^n$-torsion
points: a fixed finite number of components of dimension $b$
has at most a fixed multiple of $\ell^{2nb}$ such points,
whereas $B$ of dimension $a$ has $\ell^{2na}$. Hence $B\subset Q$.

Every component of $P_{\mathrm{red}}$ contains a torsion point.
If the component of $x$ has order $n$, then $nx\in B(k)$.
Divisibility of $B(k)$ gives $b\in B(k)$ with $nb=nx$; thus
$x-b$ is torsion. Its finite cyclic subgroup is in $Q$, and
translation proves $P_{\mathrm{red}}\subset Q$.

Choose $r$ large enough to kill every nilpotent by its $p^r$-th
power. Relative Frobenius $F^r:P\to P^{(p^r)}$ then factors
through $P_{\mathrm{red}}^{(p^r)}$. Its kernel is finite and
therefore lies in $Q$. Frobenius on the smooth group
$P_{\mathrm{red}}$ is faithfully flat. For any $x\in P(T)$,
after an fppf base change there is $b\in P_{\mathrm{red}}(T)$
with $F^r(b)=F^r(x)$. Then $x-b$ is in the finite kernel.
Both terms lie in $Q$, so $P=Q$ as fppf sheaves and hence as
group schemes.
\end{proof}

\begin{proof}[Proof of the Picard assertion in Theorem~\ref{thm:birational}]
Put $P=\Pic^\tau_{X/k}$ and $Q=\Pic^\tau_{Y/k}$. These are
proper group schemes of finite type: the identity component of
the Picard scheme of a normal proper variety is proper, and the
torsion subgroup of its N\'eron--Severi group is finite
\cite[Theorem 5.4, pp.~37--38; Remarks 5.6, p.~39, and 6.19,
p.~59]{Kleiman}.
The equality $f_*\OO_X=\OO_Y$ persists after any base change
over $k$. The projection formula then shows that a line bundle
whose pullback is trivial was already trivial. Applying this
after fppf localization proves that $Q\to P$ is a monomorphism
of relative Picard sheaves. It is proper, hence a closed immersion.

For any finite commutative group scheme $H$, let $H^D$ be its
Cartier dual. Picard--torsor duality gives, naturally for $Z=X,Y$,
\[
 H^1_{\mathrm{fppf}}(Z,H^D)
   =\Hom(H,\Pic_{Z/k})
   =\Hom(H,\Pic^\tau_{Z/k})
\]
\cite[Proposition 6.2.1, pp.~50--51]{RaynaudPicard}.
The hypotheses include properness, flatness over $k$, and
$H^0(Z,\OO_Z)=k$ universally; all hold here. A $k$-point
rigidifies the relative assertion, and finite torsors over the
algebraically closed field are trivial. The final equality holds
because a finite group has torsion image in the component group.
Finite-flat torsor purity makes the two displayed Hom groups
equal under pullback. Lemma~\ref{lem:detect} now gives $Q=P$.
\end{proof}
